\chapter{The dynamic object model}
\label{ch:dynamic}

Every backend so far spends the reflection result on \emph{framework calls}:
\code{py::class\_<Mesh>().def("at", \ldots)}, an N-API property descriptor, a
sol2 usertype. The information is consumed at generation time, and what survives
is code for one specific runtime.

The \lang{dynamic} backend spends it on \textbf{data instead}. It writes the IR
back out as static tables, so the answer to ``what does this class have'' is
still available at run time --- to a caller that was never compiled against your
headers.

\section{What it emits}

\begin{lstlisting}[style=plain]
bindings/dynamic/auto_dynamic.h     declares <lib>::register_all()  -- idempotent
bindings/dynamic/auto_dynamic.cpp   the tables
bindings/dynamic/inspect.cpp        a registry walker, for free
bindings/dynamic/CMakeLists.txt     a static library, <lib>_dynamic
\end{lstlisting}

\code{auto\_dynamic.cpp} is the only file that includes your header. It is
ordinary C++20: the metadata is data, so nothing here needs reflection to
compile.

\section{Using it}

\begin{lstlisting}[style=cpp]
#include <rosetta/dynamic.h>     // the runtime -- ordinary C++20
#include "auto_dynamic.h"        // the generated tables (scene.h is NOT included)

scene::register_all();           // publish the tables into the registry, once

for (const rosetta::dyn::MetaClass *k : rosetta::dyn::registry().classes())
    for (std::size_t i = 0; i < k->n_fields; ++i)
        std::cout << k->qualified << "::" << k->fields[i].name << '\n';
        // ... plus .type, .doc, .readonly, .range, .choices
        //     and the .get / .set thunks
\end{lstlisting}

Note what the consumer does \emph{not} do: it never includes \code{scene.h}, and
it names no bound type. That is the property everything else in this chapter
follows from.

\section{What the tables look like}

One \code{TypeDesc} per distinct type, shared by address; one \code{MetaField} /
\code{MetaMethod} / \code{MetaCtor} array per class; one \code{MetaClass} tying
them together. Straight from generated output:

\begin{lstlisting}[style=cpp]
TypeDesc td_0{.kind = Kind::number, .spelling = "double", .integral = false};
TypeDesc td_6{.kind = Kind::vector, .spelling = "std::vector<double>",
              .element = &td_0};

const MetaField k_scene__Vec3_fields[] = {
    {.name = "x", .type = &td_0, .doc = "X component",
     .annotations = k_scene__Vec3_x_anns, .n_annotations = 2,
     .get = +[](const ObjectRef &self, const ArgList &) {
         return make_any(&td_0, static_cast<scene::Vec3 *>(self.ptr)->x);
     },
     .set = +[](const ObjectRef &self, const ArgList &a) {
         static_cast<scene::Vec3 *>(self.ptr)->x = value_cast<double>(a[0]);
         return Any::none();
     }},
\end{lstlisting}

Three properties of that shape are deliberate:

\begin{itemize}
  \item \textbf{Aggregate initialisers of trivially-constructible data} ---
    \code{const char *}, pointer-and-count spans, function pointers. No
    \code{std::string}, no \code{std::vector}, no dynamic initialisation. A
    500-class library's metadata lands in \code{.rodata}, not in a static
    constructor that runs at load.
  \item \textbf{One captureless lambda per member}, decaying to a plain function
    pointer:
    \code{using Thunk = Any (*)(const ObjectRef \&self, const ArgList \&args)}.
    That is the \emph{only} callable shape in the model --- a field getter, a
    field setter, an instance method, a static method and a free function all
    reduce to it. \code{self} is the receiver, default-constructed when there
    is not one.
  \item \textbf{\code{TypeDesc} is structural, not an identity.}
    \code{std::type\_index} answers ``are these the same type''; a UI and a
    marshaller need ``what is \emph{inside} it'', so \code{element} recurses and
    \code{vector<vector<Foo>>} is describable.
\end{itemize}

\section{Why \code{register\_all()} exists}

Two classes can name each other's types, so \code{TypeDesc::cls} cannot be
filled during static initialisation without an ordering you have no way to
guarantee. The generated function adds every class, enum and function to the
registry and then calls \code{link()} to resolve those back-pointers:

\begin{lstlisting}[style=cpp]
void register_all(rosetta::dyn::Registry &r) {
    r.add_class(&k_scene__Vec3);
    r.add_class(&k_scene__Mesh);
    r.link();
}
\end{lstlisting}

It is idempotent, and it is the one line a consumer must run before querying
anything.

\section{Calling by name}

\code{Any} is the value type crossing the boundary, and its \textbf{canonical
representation} is what makes a single generic wrapper possible: whatever the
C++ said, the box holds exactly one of \code{bool}, \code{long long},
\code{double}, \code{std::string}, \code{std::vector<Any>} or an
\code{ObjectRef}.

So a caller reads a number with \code{as\_number()} without caring whether the
signature said \code{int}, \code{size\_t} or \code{uint32\_t} --- and changing
one for the other in your library is not a break for the wrapper.

\code{resolve()} scores every candidate of a name against the actual arguments
on a deliberately small ladder --- \code{exact} / \code{promote} /
\code{convert} / \code{none} --- rather than reproducing C++'s conversion
lattice. The question from a dynamically typed host is ``did the user mean this
overload'', not ``which standard conversion sequence is shorter''.

An equal best score is reported as \emph{ambiguous} rather than silently picked,
and a failure names every candidate and why it lost:

\begin{lstlisting}[style=plain]
Mesh::at(std::string) matched no overload of 2:
  at(int) -- argument 1 ("nope") is not a int
  at(int, int) -- takes 2 argument(s), 1 given
\end{lstlisting}

\section{Two consequences worth knowing}

\subsection*{What cannot be marshalled is described, not deleted}

A member whose type no gate claimed keeps its entry, with a null thunk and a
stated reason:

\begin{lstlisting}[style=cpp]
.skip_reason = "callback: a std::function parameter needs a foreign-callable adapter"
\end{lstlisting}

So a user interface can grey it out instead of the member silently vanishing.
This is the same information \code{coverage.json} carries, available at run time
rather than at review time.

\subsection*{Ownership is explicit}

\code{ObjectRef::owner} is non-null when the reference owns the object. A getter
returning \code{T\&} can hand back a \emph{borrowed} handle carrying the
parent's owner --- which pins the parent and closes the dangling-sub-object hole
that a raw-pointer return has.

\section{What it buys}

Relative to the language backends:

\begin{itemize}
  \item \textbf{Overloads come back.} The metadata keeps the whole overload set
    and scores it against the actual arguments, so the name-keyed targets'
    first-only limit (Chapter~\ref{ch:coverage}) does not apply --- and a failed
    call can name every candidate it rejected, with the reason.
  \item \textbf{UI by query, not by codegen.} A property inspector is a walk
    over \code{registry()}: \code{label} to caption, \code{doc} to tooltip,
    \code{range} to slider bounds, \code{readonly} to a disabled row,
    \code{combobox} and enumerators to a drop-down, \code{button} to an action
    row. Annotations are enforced once, in the core, instead of per backend.
  \item \textbf{Lifetime is explicit}, as above.
\end{itemize}

\section{What it costs}

Every access is a linear name lookup plus a scoring pass.

That is the right price for menus, property panels, commands and scripting
glue. It is the wrong price for bulk data --- pulling a mesh across boxes every
coordinate into an \code{Any}.

\begin{note}
Cache the resolved \code{MetaMethod}: the pointer is stable for the process.
And keep an expanded binding for hot paths --- the two projections are not
exclusive, and a library that emits both gets the property sheet from one and
the throughput from the other.
\end{note}

\section{When to reach for it}

\begin{center}
\small
\begin{tabular}{@{}L{6.4cm}L{7.0cm}@{}}
\toprule
\textbf{You want\ldots} & \textbf{Use} \\
\midrule
Python/JS/Lua callers of your API & a language backend \\
a property panel over your types & \lang{dynamic} + a walker, or the
  \lang{qt} / \lang{qml} / \lang{imgui} backends \\
a command console, a menu system, an undo stack keyed by member name &
  \lang{dynamic} \\
one scripting wrapper that works for \emph{any} library you have bound &
  \lang{dynamic} + the \code{scriptable} preset (Chapter~\ref{ch:scriptable}) \\
maximum throughput on a hot call & a language backend \\
\bottomrule
\end{tabular}
\end{center}

The worked example is \code{examples/dynamic}: one set of generated metadata
driving a terminal interpreter \emph{and} a Qt viewer --- 3-D view, property
panel, console --- neither of which names a single bound type.
